\section*{NeurIPS Paper Checklist}

\begin{enumerate}

\item {\bf Claims}
    \item[] Question: Do the main claims made in the abstract and introduction accurately reflect the paper's contributions and scope?
    \item[] Answer: \answerYes{}
    \item[] Justification: The abstract and introduction state both findings within their identified scope: the coded omission gap is indistinguishable from the label-free artifact, and within-release composition tracks standing. Section~\ref{sec:reach} reports the reliability of the coding behind them.


\item {\bf Limitations}
    \item[] Question: Does the paper discuss the limitations of the work performed by the authors?
    \item[] Answer: \answerYes{}
    \item[] Justification: Limitations appear where they bear on the argument rather than only at the end. Section~\ref{sec:data-model} bounds the estimand and states that no estimate of it is reported; Section~\ref{sec:core-boundary} states that the conditioning is a decomposition rather than a control strategy, and what a study holding disclosure labels would additionally have to argue; Section~\ref{sec:core-recentring} reports the corrections that fail and the coverage the balanced measure loses; Section~\ref{sec:reach} reports the ceiling on the design that narrows endowment uncertainty furthest, its 16 realised drops, and the underpowered null they return.


\item {\bf Theory assumptions and proofs}
    \item[] Question: For each theoretical result, does the paper provide the full set of assumptions and a complete (and correct) proof?
    \item[] Answer: \answerYes{}
    \item[] Justification: Section~\ref{sec:data-model} states one identification result, that a temporal eligibility gate cannot satisfy the exclusion restriction the design needs, and gives its argument in full: the three date variables are exactly collinear, so excluding one requires that at least two of them not move standing. Section~\ref{sec:core} shows all three move it. Appendix~\ref{sec:supporting} states the three sign restrictions used to bound the unidentified component and derives the bound from them. Everything else in Section~\ref{sec:data-model} restates published results with citation.


    \item {\bf Experimental result reproducibility}
    \item[] Question: Does the paper fully disclose all the information needed to reproduce the main experimental results of the paper to the extent that it affects the main claims and/or conclusions of the paper (regardless of whether the code and data are provided or not)?
    \item[] Answer: \answerYes{}
    \item[] Justification: Every number in the paper is written by one command into a single machine-readable file, and the data snapshot it was computed against is pinned by a per-file SHA-256 manifest in the released code. Section~\ref{sec:data} states the release key, the eligibility rule, the window width, and which estimates impose a minimum benchmark count and which do not.



\item {\bf Open access to data and code}
    \item[] Question: Does the paper provide open access to the data and code, with sufficient instructions to faithfully reproduce the main experimental results, as described in supplemental material?
    \item[] Answer: \answerYes{}
    \item[] Justification: Code and derived data are released with instructions. The independent scores are redistributed by their publisher under CC-BY and are fetched rather than copied, then verified against the pinned manifest, so an upstream change is reported rather than silently absorbed. The external leaderboard evidence is frozen with a SHA-256 per source file, because a public leaderboard moves and the argument depends on the exact values.



\item {\bf Experimental setting/details}
    \item[] Question: Does the paper specify all the training and test details (e.g., data splits, hyperparameters, how they were chosen, type of optimizer) necessary to understand the results?
    \item[] Answer: \answerYes{}
    \item[] Justification: There is no model training. The choices that affect the estimates are stated in Section~\ref{sec:data}: the release key, the rule for collapsing scaffold variants, the 182-day comparison window, which estimates impose a minimum benchmark count and which do not, and the treatment of pairs whose benchmark vintage is unverified.


\item {\bf Experiment statistical significance}
    \item[] Question: Does the paper report error bars suitably and correctly defined or other appropriate information about the statistical significance of the experiments?
    \item[] Answer: \answerYes{}
    \item[] Justification: Section~\ref{sec:data-inference} states the inference protocol. The cluster count is small and unevenly distributed, so below twelve clusters no analytic standard error is reported, coefficients in that range use a restricted wild cluster bootstrap, release-level means use a cluster sign-flip randomization test, and no result is called significant on a cluster-t alone. Table~1 reports a standard error with every coefficient, the headline contrast carries a randomization p-value, and the permutation null is reported with its mean, its spread and the number of draws reaching the observed value.


\item {\bf Experiments compute resources}
    \item[] Question: For each experiment, does the paper provide sufficient information on the computer resources (type of compute workers, memory, time of execution) needed to reproduce the experiments?
    \item[] Answer: \answerYes{}
    \item[] Justification: All analysis runs in seconds on a laptop. There is no training, no accelerator use, and no discarded run.


\item {\bf Code of ethics}
    \item[] Question: Does the research conducted in the paper conform, in every respect, with the NeurIPS Code of Ethics \url{https://neurips.cc/public/EthicsGuidelines}?
    \item[] Answer: \answerYes{}
    \item[] Justification: The work uses published benchmark scores and public release documents, involves no human subjects and no personal data, and makes no claim about any individual.



\item {\bf Broader impacts}
    \item[] Question: Does the paper discuss both potential positive societal impacts and negative societal impacts of the work performed?
    \item[] Answer: \answerYes{}
    \item[] Justification: Section~\ref{sec:external} states the consequence for a reader of a pool-relative leaderboard, namely that such a number describes standing within one vintage of the evaluated pool rather than a property of the model. The conclusion states that whether the record needed to measure withholding should grow is a reporting-standard question. The paper does not attribute concealment to any named provider, because the design does not support that.


\item {\bf Safeguards}
    \item[] Question: Does the paper describe safeguards that have been put in place for responsible release of data or models that have a high risk for misuse (e.g., pre-trained language models, image generators, or scraped datasets)?
    \item[] Answer: \answerNA{}
    \item[] Justification: No model, dataset or other asset with misuse potential is released. The release is analysis code and derived tables built from already public sources.


\item {\bf Licenses for existing assets}
    \item[] Question: Are the creators or original owners of assets (e.g., code, data, models), used in the paper, properly credited and are the license and terms of use explicitly mentioned and properly respected?
    \item[] Answer: \answerYes{}
    \item[] Justification: The independent scores are credited to their publisher and used under CC-BY, and the external leaderboard is cited to its published description. Both are named in the text and in the released code.


\item {\bf New assets}
    \item[] Question: Are new assets introduced in the paper well documented and is the documentation provided alongside the assets?
    \item[] Answer: \answerYes{}
    \item[] Justification: The released code documents how each number is produced, the snapshot manifest records the exact data build, and the frozen external evidence records a source URL, a retrieval date and a SHA-256 for every file.


\item {\bf Crowdsourcing and research with human subjects}
    \item[] Question: For crowdsourcing experiments and research with human subjects, does the paper include the full text of instructions given to participants and screenshots, if applicable, as well as details about compensation (if any)?
    \item[] Answer: \answerNA{}
    \item[] Justification: The work involves no crowdsourcing and no human subjects.


\item {\bf Institutional review board (IRB) approvals or equivalent for research with human subjects}
    \item[] Question: Does the paper describe potential risks incurred by study participants, whether such risks were disclosed to the subjects, and whether Institutional Review Board (IRB) approvals (or an equivalent approval/review based on the requirements of your country or institution) were obtained?
    \item[] Answer: \answerNA{}
    \item[] Justification: The work involves no human subjects, so no review board approval was required.


\item {\bf Declaration of LLM usage}
    \item[] Question: Does the paper describe the usage of LLMs if it is an important, original, or non-standard component of the core methods in this research? Note that if the LLM is used only for writing, editing, or formatting purposes and does \emph{not} impact the core methodology, scientific rigor, or originality of the research, declaration is not required.
    %this research?
    \item[] Answer: \answerYes{}
    \item[] Justification: An earlier pass used a large language model to assign the disclosure-coding categories; those labels were cleared and none enters. The categories now standing were hand-assigned by a human coder from archived release-time documents under the frozen protocol, with an independent second coder blind-coding a provider-stratified 20\% sample (Section~\ref{sec:reach}), so no number reported here depends on a model's judgment. Every number is produced by the released code from the pinned snapshot. Language models were otherwise used only for writing and code assistance.


\end{enumerate}
